\chapter{The Empirical and Logical Traditions}
\label{ch:empiricism}

\begin{chapterguide}
Logical empiricism is reconstructed as a family of projects rather than a
single doctrine. The chapter then examines its achievements and difficulties:
the demand for public evidence, formal clarification, reduction of theoretical
claims, probability, confirmation, and the problems of meaning, holism, and
theory-ladenness. It closes with the lesson that empiricism supplies a
discipline of contact but not a complete account of objectivity.
\end{chapterguide}

The label \emph{logical positivism}\index{logical positivism} is often used as
if it named a uniform creed. Historically, the movement was more varied. The
Vienna Circle, the Berlin group, and associated philosophers disagreed about
probability, protocol statements, physicalism, conventionality, and the proper
role of metaphysics. The later label \emph{logical empiricism} is often more
accurate because several figures abandoned the strongest version of a
verification principle and developed sophisticated accounts of theoretical
terms, confirmation, and scientific language.

The broad ambition was clear: make scientific claims answerable to experience
and make their logical structure explicit. This ambition remains part of any
defensible epistemology. The problem is not the desire for evidence. The
problem is the idea that evidence can be specified without theoretical,
instrumental, historical, or social conditions.

\section{Meaning, evidence, and the rejection of pseudo-problems}

Logical empiricists sought a criterion for distinguishing statements that are
empirically meaningful from those that merely appear to make claims. A
mathematical statement could be meaningful through formal relations. An
empirical statement could be meaningful through its connection to possible
experience. Traditional metaphysical sentences were criticized when they
seemed to lack a method of verification or a clear inferential role.

The strongest charitable interpretation is methodological rather than
dismissive. The empiricists were not simply saying that everything not
immediately observable is nonsense. They were asking what difference a claim
makes to possible observation, calculation, explanation, or action. A statement
about an electron could be meaningful if it entered laws and experimental
procedures that generated testable consequences. A statement about a
transcendent entity that made no difference to any possible inquiry would not
belong to empirical science.

This motivation can be retained without adopting a strict verification
criterion. Scientific meaning is partly a matter of inferential role: what
would the claim predict, explain, license, rule out, or change in a measurement
procedure? The fact that a term is theoretical does not disqualify it. It
requires a map of how the term connects to observation and intervention.

\section{The appeal of protocol statements}

A basic problem in empiricism is the status of the observational basis.
Scientific claims are tested against observations, but observations are
reported in language and obtained under conditions. Is there a neutral layer of
protocol statements that is independent of theory?

Different empiricists proposed different answers. Some emphasized immediate
experience; others moved toward physical language, public records, or
intersubjective agreement. The movement away from private sense data was a
major strength. A scientific observation must be inspectable by others, linked
to a procedure, and embedded in a record. A private sensation can be evidence
for a person's experience, but it is not yet a public measurement.

Yet public accessibility is not the same as theoretical neutrality. To report
that a pointer moved, one needs a concept of a pointer and a rule for reading
it. To report a spectral line, one needs an instrument, a calibration procedure,
and a model connecting detector counts with wavelength. Observation sentences
are not arbitrary, but they are not self-authenticating either.

\begin{principlebox}
\textbf{Empirical contact principle.} A scientific observation is not a bare
given. It is a publicly inspectable result of a procedure that connects a
physical interaction to a claim under stated assumptions.
\end{principlebox}

This principle keeps the empiricist demand for public contact while refusing
the fantasy of an uninterpreted foundation. The later chapters will call the
stated assumptions a \emph{measurement path}: the chain from target, through
instrument and calibration, to recorded datum and interpreted result.

\section{Confirmation and the problem of induction}

Suppose a scientist observes many white swans. The observations support the
claim that all swans are white, but they do not entail it. The next swan could
be black. This is the familiar problem of induction. The problem is not
specific to swans. A theory may fit every observation collected so far and
still fail tomorrow.

Logical empiricists responded in several ways. Some treated confirmation as a
degree of support rather than proof. Reichenbach developed a probabilistic
framework in which inductive rules were justified pragmatically: if a stable
frequency existed, a suitable rule would converge toward it. Carnap attempted
to formalize degrees of confirmation. Others accepted that scientific inference
is ampliative and emphasized the practical success of prediction.

The important lesson is that evidence should change belief without pretending
to establish certainty. A rational method needs at least three ingredients:

\begin{enumerate}
\item a set of candidate claims;
\item information about how likely the evidence would be under those claims;
\item a rule for representing uncertainty and updating.
\end{enumerate}

Bayesian epistemology provides one such rule. If \(H\) is a hypothesis and
\(E\) is evidence, then Bayes' theorem is

\[
  \Prob(H\mid E)
  =
  \frac{\Prob(E\mid H)\Prob(H)}{\Prob(E)}.
\]

Here \(\Prob(H)\) is the prior probability assigned before \(E\), \(\Prob(E\mid
H)\) is the likelihood of observing \(E\) if \(H\) is true, and \(\Prob(H\mid
E)\) is the posterior probability after the evidence. The denominator
\(\Prob(E)\) makes the posterior probabilities sum to one across the relevant
alternatives.

The equation is a consistency rule, not a machine that discovers the correct
prior. If two investigators assign different priors, they may initially
disagree even after observing the same evidence. They may nevertheless converge
when the likelihood ratio is sufficiently strong. This observation will matter
later: public objectivity cannot be identified with identical starting
beliefs, but it can be strengthened by evidence that sharply discriminates
among alternatives.

\section{Theoretical terms are not optional decorations}

A strict empiricism sometimes imagines that science could be translated into
a language containing only observable predicates. But theoretical terms do
more than abbreviate observations. They organize relations among observations,
generate new questions, and make interventions possible.

Consider temperature. A thermometer reading is not identical to temperature.
The reading is produced by a physical process whose response is calibrated
against a scale. The theoretical concept of temperature links different
instruments, equilibrium conditions, equations, and predictions. Removing the
term would not make the science more empirical; it would make the connections
harder to express.

The same point applies to genes, fields, utility, inflation expectations, and
neural representations. Some of these terms refer to entities, some to
properties, some to model variables, and some to institutional constructs.
They should not all receive the same ontological status. But they cannot be
evaluated merely by asking whether they are directly observed.

A scientific term earns epistemic standing by the role it plays in a network
of constraints. It should connect to measurement, support counterfactuals,
survive alternative operationalizations, or make possible reliable
intervention. The term becomes suspicious when it is insulated from every
possible correction and is used only to redescribe its own successes.

\section{Holism and the Duhem--Quine problem}

An isolated hypothesis rarely entails an observation by itself. A prediction
typically depends on auxiliary assumptions about initial conditions,
instruments, background laws, numerical approximations, and data processing.
When the prediction fails, several components may be responsible. This is the
problem associated with Duhem and later Quine.

Quine's critique of the analytic--synthetic distinction extended the point:
beliefs form a web, and experience bears on the web as a whole rather than on
single statements one at a time \citep{duhem1906,quine1951}. The strongest
lesson is epistemic humility. A failed prediction does not automatically
identify which sentence is false.

The bad lesson would be that no claim can be tested. Scientific practice
responds by designing tests that vary or stabilize different parts of the
network. Independent instruments can challenge calibration assumptions.
Different initial conditions can test a law across regimes. Replication teams
can use the same data with alternative pipelines. Interventions can isolate
causal components. A theory can therefore be tested holistically, but not
indiscriminately.

\begin{table}[ht]
\centering
\smalltable
\begin{tabularx}{\textwidth}{p{0.29\textwidth}YY}
\toprule
Empiricist insight & What it protects & What it cannot by itself supply\\
\midrule
Public evidence & Shared access and correction & A theory-neutral observation language\\
Formal clarity & Explicit inferential commitments & A choice of relevant variables\\
Probability & Graded belief and uncertainty & A unique prior or causal interpretation\\
Operational connection & Contact between concepts and procedures & A guarantee that the operation measures the intended target\\
\bottomrule
\end{tabularx}
\caption{What the logical empiricist tradition contributes, and where it needs
supplementation.}
\label{tab:empiricist-lessons}
\end{table}

\section{What survives}

The logical empiricist tradition contributes four durable norms.

First, scientific claims should be connected to public procedures. A claim that
cannot alter any observation, intervention, or inferential practice is not a
scientific claim merely because it uses technical words.

Second, the relation between theory and evidence should be made explicit.
Scientists should state what assumptions a result depends on and what would
change if those assumptions were relaxed.

Third, uncertainty should be represented rather than hidden behind categorical
language. A point estimate without an error model can be less informative than
a range with clear limitations.

Fourth, philosophical arguments benefit from formal reconstruction. Logic and
probability do not replace empirical inquiry, but they expose ambiguities and
show exactly where an inference requires an extra premise.

What does not survive is the idea of a purely given foundation, a universal
translation into observation language, or an automatic rule for induction.
Objectivity will require a broader account of constraints: experimental,
statistical, causal, historical, and social.

\section*{Chapter summary}

Logical empiricism is best understood as a diverse attempt to connect meaning
and justification to public evidence and logical structure. Its strongest
achievements are the demand for inspectable procedures, explicit inferential
commitments, probabilistic uncertainty, and clarity about theoretical terms.
Its weaknesses arise when observation is treated as theory-neutral or when
confirmation is expected to yield certainty. Holism shows that tests involve
networks of assumptions, but independent instruments, interventions, and
replications can still distribute and localize criticism.

\begin{discussion}
\begin{enumerate}
\item Does a claim become scientific when it has a possible observation, or
does it also need a reliable measurement path?
\item Should a theoretical entity be accepted if it improves prediction but
does not support intervention?
\item How should two investigators with different priors resolve a disagreement?
\end{enumerate}
\end{discussion}

\begin{furtherreading}
Read Reichenbach (1938) for a mature probabilistic empiricism, Hempel (1965)
for the logic of explanation, Quine (1951) for holism, and the historical
overview of logical empiricism in the relevant essays collected by Carnap and
his contemporaries.
\end{furtherreading}

